\chapter{Conclusion\label{cha:chapter6}}

\section{Summary}

This thesis overhauled Fainder's query execution pipeline into a compiled Rust extension module. The new engine was used for an ablation study across more than twenty hardware-conscious optimisation axes. The dominant outcome of this study is structural. The engineering speedup, while large, is mechanistically secondary to it.

On the Sapphire Rapids 8468H micro-architecture, Fainder's search phase is bound by four distinct performance ceilings: 

\begin{enumerate}
    \item Single-thread compute / L1 latency (\(t\leq 8\); candidate, refuted)
    \item Shared L3 capacity (\(t\approx 16\))
    \item On-socket memory-traffic pressure (\(t\in[32, 48]\) on one socket)
    \item Cross-socket UPI interconnect (\(t > 48\))
\end{enumerate}

Ceiling (i) is the literature's default assumption for binary-search workloads, its refutation on this workload is itself a finding. The ceilings were established through Chapter~\ref{cha:chapter4}'s perf counters. Ceiling (iii) and (iv) were separated by a four-configuration \texttt{numactl} probe.

Across five independent ablations on five different axes, a pattern emerged: a less fine-grained optimisation had already addressed the binding ceiling, and the finer-grained optimisation's overhead was a net regression rather than diminishing return. Although each instance had an argument to compose positively, each was falsified at scale by hardware counter evidence. The cross-axis breadth of the five instances rules out a per-axis coincidence and establishes the pattern as structural. The five ablations were: emit-phase ID packing over \bcombo{} (§\ref{sec:exp:c3:packed}), K-batch cluster amortisation over \texttt{f16} (§\ref{sec:exp:wf:qb}), morsel-driven phase coordination over query-batch (§\ref{sec:exp:wf:morsel}), per-cluster reindexing over packed-ids (§\ref{sec:exp:reindex}), and page-granular NUMA interleave over first-touch placement (§\ref{sec:exp:c4}). 

This pattern makes \textbf{dispatch-on-regime} the only near-optimal deployment strategy, which selects build features and deployment flags per workload regime. This dispatch policy was validated 17/18 exactly and 18/18 within 5\% in-distribution, and 6/6 within 5\% out-of-distribution, all six within 1\%, for a combined 23/24.

The engineering results are secondary to the structural finding: Rust achieves a $530$--$598\times$ end-to-end speedup over the Python baseline at the peak-scaling cell ($t=16$) across three independent dataset scales. This end-to-end ratio is dominated by Rust avoiding Python's result-dict construction (§\ref{sec:exp:e2e:mat}). The search phase itself scales from $22.7\times$ on the smallest calibration dataset to $1.34\times$ on the largest, bending the curve toward the shared hardware ceilings. All these speedups preserve precision $=$ recall $= 1.000$. The speedup depends on the regime, since negative composability makes the best composition for one regime, regress at another.

\section{Contributions}

The thesis makes five primary contributions.

\begin{enumerate}
    \item \textbf{Engineering re-implementation.} Rust port of Fainder's query execution pipeline, validated at precision $=$ recall $= 1.000$. Integrated into the existing Python API via PyO3, delivering a speedup across three independent dataset scales.
    
    \item \textbf{Four-ceiling characterisation.} Identification of four performance ceilings mentioned above, on the Sapphire Rapids 8468H micro-architecture, and the regimes under which each ceiling binds. 

    \item \textbf{Negative-composability pattern.} Five independent instances of composability where a second optimisation regresses at scale. An ablation that would compose positively would falsify this pattern.

    \item \textbf{Pre-flight ceiling-identification.} A probe template diagnoses the binding ceiling at the \emph{target composition} before committing to implementation. When multiple stacked ceilings exist, the ablation methodology (implement, measure, accept) is methodologically insufficient.
    
    \item \textbf{Dispatch-on-regime.} Automatic policy that selects build features and deployment flags per workload regime. Valid even across out-of-distribution scenarios.
\end{enumerate}

Contributions (1) and (5) are engineering artefacts, while (2), (3), and (4) are scientific and methodological findings.

\section{Limitations}

\paragraph{Single micro-architecture.}
All experiments run on a single Sapphire Rapids 8468H configuration. The ceilings introduced in this thesis are based on this micro-architecture. A different micro-architecture (Genoa, Granite Rapids, Grace) would most likely have different per-ceiling thresholds. The methodological contribution (pre-flight ceiling-identification) generalises across micro-architectures, the specific ceiling thresholds and dispatch policy entries do not.

\paragraph{Single-machine evaluation.}
All experiments run on shared memory. This might introduce a fifth ceiling of Network latency, in case of sharing the index across multiple machines. Asynchronous I/O and work offloading become candidate optimisations in that case.

\paragraph{Density-range scope.}
The dispatch-on-regime policy is validated at \(n_\text{hists}/n_\text{clusters}\) densities of $\sim$8,200 (c1024\_56gb) and $\sim$26,300 (c256\_56gb). The policy is not validated at densities below $\sim$8,200 (c4096\_56gb) or above $\sim$26,300 (c128\_56gb), both outside the validated range. The policy's density axis is therefore a prediction rather than a validated entry.

\section{Future Work}

Future work includes three experiments that would extend the structural findings.

\paragraph{(F1) Transparent Huge Pages probe across all builds.}
The §\ref{sec:exp:reindex} mechanism localised the local-ids regression. The bw=13 decoder's access pattern over the packed buffer drives page-walk pressure. With Transparent Huge Pages enabled ($2$\,MiB pages instead of $4$\,KiB), TLB-miss frequency drops by $512\times$ per covered page, and the page-walker's $\sim\!100$-cycle penalty per miss correspondingly shrinks.

\paragraph{(F2) NUMA-aware build for $t > 48$.}
§\ref{sec:exp:c4} identified cross-socket UPI as the fourth ceiling and showed that single-socket placement saves 11--13\% at $t \leq 48$. Single-socket placement is infeasible at $t >  48$ due to socket 0 having only 48 physical cores. A NUMA-aware build would shard the cluster set across the two sockets at index-load time and pin Rayon workers to physical cores within their NUMA node. Queries would then iterate over node-local clusters preferentially.

\paragraph{(F3) Smaller follow-ons.}

\textit{Adaptive precision selection.} Optimal f16 vs f32 precision depends on the regime. F32 wins at low $t$, while f16 wins at high $t$. An adaptive scheme at runtime that selects precision based on the thread count (or L3-miss rate) would provide the best of both.

\textit{Column-centric engine for small datasets.} Column-centric outperformed row-centric at $t=1$ by $2.17$--$2.63\times$ (§\ref{sec:exp:wf:columnar}). Broadening its win region would require addressing the cross-cluster load imbalance that causes columnar to lose by a factor of 2 at $t \geq 8$.

\section{Closing Remark}

The engineering speedups produced in this thesis are specific to this codebase, micro-architecture, and workload. The structural finding is what generalises. On hardware-near workloads the cost of regression for optimisation $B$ on an already applied optimisation $A$ depends on which ceiling has been collected already, making it invisible to per-operator cost models. This makes our finding clear: identify the binding ceiling at the target composition before committing the implementation cost of the next ablation.
